%%%%%%%%%%%%%%%%%%%%%%%%%%%%%%%%%%%%%%%%%%%%%%%%%%%%%%%%%%%%%%%%%%%%%%%%%%%%%%%%
%% 
\cleardoubleoddpage%  Make sure to start each chapter on a new odd page
\chapter{Conclusion and Outlook}
\label{chapter:conclusion}
In conclusion, we started by introducing several notions such as the scheduling problem we designed the \ac{IR} to solve. Following this, we defined the language design, including several iterations and decisions we made along the way. Next, we showed how our language is designed and why we made certain decisions, indicating the descriptive power of our language. After defining the language, we formalized the constraints mathematically, giving a clear definition of what a solution is and how to compare different solutions. In the process of defining the formal framework, we show required properties, which pose as input for the solver. Right after we defined the formal definition of a valid solution. Afterwards, we defined how to compare different solutions using cost functions, which can be used to find optimal solutions. Finally, we introduced some notions of solvers to get an idea of what the language would be best used for, given that the issue of different formulations differing in solver efficiency persists for \ac{ILP} \cite{bib:1994-chaudhuri}. \\
Overall we can state, that we successfully defined an \Ac{IR} for scheduling problems. This language is designed to pose as an intermediate translation step between any high-level language and low-level solver-ready constraints. Both the high and low-level languages are replaceable and not bound to any specific implementation. This allows for a wide range of applications, as long as the high-level language can be transpiled into our \Ac{IR} and the low-level solver can handle the constraints we defined. \\
For future work, we see several possibilities. First and foremost, implementing a solver which can handle the defined constrains is of utmost importance. The defined formal framework can be used as a basis for such an implementation. Such an implementation - or rather the output of one - can be checked with an existing implementation of a solution verifier \footnote{\href{https://github.com/h0p3zZ/BScThesis/tree/dev/_prog}{Repository with the implementation of a syntax compiler for the \Ac{IR} and a solution verifier.}}. Furthermore, benchmarking different types of solver with our defined constraints would be an interesting avenue for future research. Additionally, implementing a transpiler from a high-level language into our \Ac{IR} would be beneficial to showcase the usability of our language. Finally, one also has to give an intuitive way of defining tasks in a high-level language. Most likely such an input, would be in a natural language, which could be parsed using \ac{NLP} and/or \ac{LLM} techniques.
%% 
%%%%%%%%%%%%%%%%%%%%%%%%%%%%%%%%%%%%%%%%%%%%%%%%%%%%%%%%%%%%%%%%%%%%%%%%%%%%%%%%
